% Title etc.
\title{A FAIR Comic about Research Data Infrastructure}
\subtitle{Part 1: From Charlemagne to ing.grid}
\shorttitle{A FAIR Comic}

% Page 1
\FDISetText{1}{1}{1}{From Charlemagne to Little Charles.\\ From Carolingian minuscule via Middle Francia to \hspace*{0.5cm} NFDI4ING.*\\Research data infrastructure is and remains an exciting subject.}
\FDISetText{1}{1}{2}{West Francia}
\FDISetText{1}{1}{3}{Middle Francia}
\FDISetText{1}{1}{4}{East Francia}
\FDISetText{1}{1}{5}{North Sea}
\FDISetText{1}{1}{6}{Baltic\\Sea}
\FDISetText{1}{1}{7}{Kingdom\\of\\England}
\FDISetText{1}{1}{8}{Ireland}
\FDISetText{1}{1}{9}{Atlantic\\Ocean}
\FDISetText{1}{1}{10}{Kingdom\\of Asturias}
\FDISetText{1}{1}{11}{Brussels}
\FDISetText{1}{1}{12}{Aachen}
\FDISetText{1}{1}{13}{Luxem\FDIStyledText{8pt}{2pt}{-}\\[-2pt]bourg}
\FDISetText{1}{1}{14}{Verdun}
\FDISetText{1}{1}{15}{Darmstadt}
\FDISetText{1}{1}{16}{Strasbourg}
\FDISetText{1}{1}{17}{Rome}
\FDISetText{1}{1}{18}{Emirate of Córdoba}
\FDISetText{1}{1}{19}{Kingdom of Fez}
\FDISetText{1}{1}{20}{Mediterranean Sea}
\FDISetText{1}{1}{21}{Benevento}
\FDISetText{1}{1}{22}{Serbia}
\FDISetText{1}{1}{23}{Slavic peoples}
\FDISetText{1}{1}{24}{* NFDI4ING is the consortium for engineering sciences within the German National Research Data Infrastructure (NFDI)}
\FDISetText{1}{1}{25}{Limes Sorabicus}


% Page 2
\FDISetText{2}{1}{1}{In the 8th/9th century, Charlemagne began building a sustainable knowledge infrastructure through his centralized culture of writing. He relied on a network of scribes\\in monasteries and at his court in Aachen.}
\FDISetText{2}{2}{1}{To know another\\language is like possessing\\a second soul.}
\FDISetText{2}{3}{1}{Charlemagne could not write.\\When signing documents, he merely added the so-called completion stroke to the Karolus monogram.}
\FDISetText{2}{4}{1}{However, Charlemagne promoted education by having scholars dictate, copy, and distribute documents.}
\FDISetText{2}{4}{2}{Books out!\\Class test!}
\FDISetText{2}{5}{1}{To ensure the readability and consistency of documents throughout his vast empire, he commissioned the development of the \enquote{Carolingian minuscule}.}
\FDISetText{2}{5}{2}{Only lowercase letters (minuscule)}
\FDISetText{2}{5}{3}{Even, open, and calm typeface}
\FDISetText{2}{5}{4}{The \enquote{i} was still written without a dot}
\FDISetText{2}{5}{5}{Four-line system for better readability}
\FDISetText{2}{5}{6}{charles is dumb}
\FDISetText{2}{6}{1}{The \enquote{s} was written\\until the 12th century\\as a \enquote{long s}\\character}
\FDISetText{2}{6}{2}{The letter \enquote{j}\\did not yet exist\\in the 8th or\\9th century}
\FDISetText{2}{6}{3}{Unconnected letters for better readability}
\FDISetText{2}{6}{4}{\FDIUseText{2}{5}{6}}
\FDISetText{2}{6}{5}{that is iust\\not true}

% Page 3
\FDISetText{3}{1}{1}{Charlemagne's sons and grandsons continued this analogue research infrastructure based on parchment and ink.}
\FDISetText{3}{1}{2}{Charlemagne's\\clan}
\FDISetText{3}{1}{3}{(A selection)}
\FDISetText{3}{1}{4}{Charlemagne}
\FDISetText{3}{1}{5}{Pepin the\\[-2.5pt]Hunch-\\[-2.5pt]back}
\FDISetText{3}{1}{6}{Charles\\[-2.5pt]the\\[-2.5pt]Younger}
\FDISetText{3}{1}{7}{Carloman}
\FDISetText{3}{1}{8}{Louis\\[-2.5pt]the\\[-2.5pt]Pious}
\FDISetText{3}{1}{9}{Lothair I}
\FDISetText{3}{1}{10}{Pepin}
\FDISetText{3}{1}{11}{Louis II\\[-2.5pt]the\\[-2.5pt]German}
\FDISetText{3}{1}{12}{Charles II\\[-2.5pt]the\\[-2.5pt]Bald}
\FDISetText{3}{2}{1}{One of his grandsons, \FDIUseText{3}{1}{9}, played a\\major role in the Treaty\\of Verdun in 843.}
\FDISetText{3}{3}{1}{This treaty is regarded as the birth of the predecessors of France and Germany (\FDIUseText{1}{1}{2} and \FDIUseText{1}{1}{4}).}
\FDISetText{3}{3}{2}{\FDIUseText{1}{1}{2}}
\FDISetText{3}{3}{3}{\FDIUseText{1}{1}{4}}
\FDISetText{3}{4}{1}{Lothair himself\\received\\Middle Francia.\\It stretched\\as a long,\\narrow corridor\\from the\\North Sea to the\\Mediterranean.}
\FDISetText{3}{4}{2}{\FDIUseText{1}{1}{3}}
\FDISetText{3}{5}{1}{This \FDIUseText{1}{1}{3} became the nucleus\\of European integration.}
\FDISetText{3}{6}{1}{Many institutions of the\\European Union are now located along\\this former corridor\\(Brussels, Luxembourg, Strasbourg, Rome).}

% Page 4
\FDISetText{4}{1}{1}{And that is how the Comiciade\\at the Technology Centre on Europaplatz\\in Aachen opens its doors to comic fans from\\all over the world.}
\FDISetText{4}{1}{2}{Are all the\\speech bubbles in\\the right\\place?}
\FDISetText{4}{1}{3}{This way\\to the DEUS\\AIX MACHINA\\exhibition!}
\FDISetText{4}{2}{1}{Preparations are already under way\\at RWTH Aachen's C.A.R.L. as well.}
\FDISetText{4}{2}{2}{The trend\\is increasingly\\moving towards\\data storytelling!}
\FDISetText{4}{2}{3}{Has\\Little Charles\\already \\arrived?}
\FDISetText{4}{3}{1}{Mmmmhh....\\Not bad...}
\FDISetText{4}{4}{1}{Knock}
\FDISetText{4}{5}{1}{Hello,\\Évariste!}
\FDISetText{4}{6}{1}{Ah,\\there you\\are,\\Alfred!}

% Page 5
\FDISetText{5}{1}{1}{I really enjoyed your comic\\about your time at the Chair~of~Methodology\\of Artificial Intelligence. What do you think\\about making the world's first FAIR comic now?}
\FDISetText{5}{2}{1}{The world's\\first FAIR comic\\ever?}
\FDISetText{5}{3}{1}{I mean FAIR in the context of\\research data infrastructure (RDI).}
\FDISetText{5}{3}{2}{The FAIR principles are:\\Findable, Accessible,\\Interoperable, and\\Re-usable.*}
\FDISetText{5}{3}{3}{*Findable, accessible, interoperable, and re-usable}
\FDISetText{5}{4}{1}{You'll have to\\explain that to me,\\Évariste.}
\FDISetText{5}{5}{1}{It's quite simple. On the back\\of your comic, you included an ISBN\\so it can be found more easily\\in bookshops.}
\FDISetText{5}{6}{1}{Yes, the ISBN reveals the\\product, language area, publisher\\number, and the title\\number.}
\FDISetText{5}{6}{2}{Product\\(book)}
\FDISetText{5}{6}{3}{Publisher number\\(Granus Verlag)}
\FDISetText{5}{6}{4}{Title\\number}
\FDISetText{5}{6}{5}{978-3-9822083-4-3}
\FDISetText{5}{6}{6}{Country identifier (Germany,\\Austria, and Switzerland)}
\FDISetText{5}{6}{7}{Check\\digit}

% Page 6
\FDISetText{6}{1}{1}{You need an ISBN to be included in\\the German Books in Print directory (VLB).\\If you are not listed there, it is very difficult\\to get your book into shops.}
\FDISetText{6}{2}{1}{The VLB is the\\central, independent\\database and research tool\\for the German-language\\book trade.}
\FDISetText{6}{3}{1}{Interesting. We researchers\\use a DOI (Digital Object Identifier)\\instead of an ISBN.}
\FDISetText{6}{4}{1}{It is a unique, persistent identifier for objects of all kinds. An object may have several DOIs, but a DOI always points to only one object.}
\FDISetText{6}{4}{2}{DOI resolver}
\FDISetText{6}{4}{3}{Repository}
\FDISetText{6}{4}{4}{https://doi.org/10.5281/zenodo.19468332}
\FDISetText{6}{4}{5}{Prefix}
\FDISetText{6}{4}{6}{Object}
\FDISetText{6}{4}{7}{Or as a\\QR code}
\FDISetText{6}{5}{1}{A partner of\\\hspace*{0.75cm} NFDI4ING, the DataCite association,\\assigns DOIs.*}
\FDISetText{6}{5}{2}{*DataCite is a registration agency for DOIs}
\FDISetText{6}{6}{1}{DOIs are the standard in academic publications for precise source references.}
\FDISetText{6}{6}{2}{And how\\exactly do\\they\\work?}

% Page 7
\FDISetText{7}{1}{1}{They point to\\digital objects.}
\FDISetText{7}{1}{2}{Digital object 1}
\FDISetText{7}{1}{3}{Digital object 2}
\FDISetText{7}{1}{4}{Digital object 3}
\FDISetText{7}{1}{5}{Digital object 4}
\FDISetText{7}{1}{6}{Digital object 5}
\FDISetText{7}{1}{7}{Digital object 6}
\FDISetText{7}{2}{1}{Cool. This means academic\\works remain permanently citable without\\links going “dead,” right?}
\FDISetText{7}{3}{1}{Exactly! Every digital object,\\whether an article, data, or a video, receives\\a distinctive code.}
\FDISetText{7}{4}{1}{Because the DOI cannot\\be changed, it serves as a\\reliable reference.}
\FDISetText{7}{5}{1}{Hmmm ...\\What is the best\\way to\\illustrate\\that?}
\FDISetText{7}{6}{1}{Quite simple.\\We researchers\\like to use the\\tin-can\\metaphor.}

% Page 8
\FDISetText{8}{1}{1}{Beans}
\FDISetText{8}{1}{2}{And this is\\how the comparison between a tin can and\\a digital object works:\\Just as the\\contents of the\\tin can are described\\on its label, the data content of the\\Digital Object (DO)\\is described.}
\FDISetText{8}{2}{1}{This diagram illustrates very\\clearly how a DO is structured.}
\FDISetText{8}{2}{2}{Bit sequence}
\FDISetText{8}{2}{3}{Metadata\\ \& \\context}
\FDISetText{8}{2}{4}{Operations}
\FDISetText{8}{2}{5}{Persistent Identifier}
\FDISetText{8}{3}{1}{At the centre is the bit sequence.\\It consists of data in the form of\\zeros and ones.}
\FDISetText{8}{4}{1}{In our tin-can metaphor,\\this corresponds to the contents of the\\can: the beans.}
\FDISetText{8}{5}{1}{The metadata surround it.\\They describe the contents in words.\\In our case, the label says\\ \enquote{\FDIUseText{8}{1}{1}}.}
\FDISetText{8}{5}{2}{\FDIUseText{8}{1}{1}}

% Page 9
\FDISetText{9}{1}{1}{The next circle, containing\\the operations, explains what can be\\done with the contents.}
\FDISetText{9}{1}{2}{Operations}
\FDISetText{9}{2}{1}{NFDI4ING is working to define more\\than just a download function\\within the operations.}
\FDISetText{9}{3}{1}{The outer circle is the Persistent\\Identifier (PID). PID is the umbrella term\\for a range of systems that make digital\\objects findable.}
\FDISetText{9}{3}{2}{PID}
\FDISetText{9}{3}{3}{In our case, it is a DOI.}
\FDISetText{9}{4}{1}{However, there are also other systems, such as URN, Handle, or ORCID.*}
\FDISetText{9}{4}{2}{*URN = Uniform Resource Name, Handle = Handle System; ORCID = Open Researcher and Contributor ID}
\FDISetText{9}{5}{1}{So every DOI\\is a PID, but not\\every PID is automatically\\a DOI.}
\FDISetText{9}{5}{2}{Phew!\\You're really\\getting deep into\\the canned goods,\\Évariste!}
\FDISetText{9}{6}{1}{HA HA HA!\\Yes, I'm making you\\data-literate right now,\\Alfred!}
\FDISetText{9}{6}{2}{You're\\making me\\data-literate?}

% Page 10
\FDISetText{10}{1}{1}{In RDI, the term \enquote{data literacy} describes the process of acquiring the fundamental skills and knowledge required to read and write research data.}
\FDISetText{10}{2}{1}{If nobody can read them, even\\the best archive is useless, and the\\data remain silent.}
\FDISetText{10}{2}{2}{Charlemagne already knew that. He understood that infrastructure without education is useless.}
\FDISetText{10}{3}{1}{You have to\\learn to read\\the language\\of data.}
\FDISetText{10}{4}{1}{Charlemagne placed great importance on clerics and administrative officials being able to read and write {\textemdash} also so that they could copy texts correctly. Faulty data transmission was already a problem back then.}
\FDISetText{10}{5}{1}{In other words,\\copy-and-paste errors in the 9th century.}
\FDISetText{10}{5}{2}{HA HA HA!}
\FDISetText{10}{5}{3}{Now that's\\what I call\\interoperability\\between the past\\and the present!}
\FDISetText{10}{6}{1}{What are\\you laughing\\about?}

% Page 11
\FDISetText{11}{1}{1}{Hello, Tobias. We are currently\\creating a FAIR comic, and\\it is proving very entertaining.}
\FDISetText{11}{2}{1}{That sounds exciting.\\Yesterday, I submitted my latest paper to ing.grid. You could do the same\\with this comic.}
\FDISetText{11}{2}{2}{ing.grid?\\What is\\that?}
\FDISetText{11}{3}{1}{ing.grid is a\\Diamond Open Access\\journal of NFDI4ING.}
\FDISetText{11}{3}{2}{It\\also\\has a\\preprint\\server.}
\FDISetText{11}{4}{1}{That's a great idea, Tobias.\\ If we submit it there, the comic would already be visible on the preprint server and accessible to everyone.}
\FDISetText{11}{5}{1}{Exactly! A community review would be\\the first step there. That means the\\comic could already be commented on\\and evaluated.}
\FDISetText{11}{6}{1}{Great!\\Where should I put\\my completion\\stroke?}

% Page 12
\FDISetText{12}{1}{1}{Just a moment!\\First, we have to translate the comic into the target language. And there is already\\a guide for that.}
\FDISetText{12}{2}{1}{This guide can be found\\under this HANDLE.}
\FDISetText{12}{2}{2}{http://hdl.handle.net/2268.2/12918}
\FDISetText{12}{2}{3}{The citation looks like this: \cite{simons_traduction_2021}}
\FDISetText{12}{3}{1}{Okay.\\Then let's\\upload it as a PDF now.*}
\FDISetText{12}{3}{2}{* PDF = Portable Document Format}
\FDISetText{12}{4}{1}{Aaaand...\\Here we go!}
\FDISetText{12}{4}{2}{CLICK}
\FDISetText{12}{5}{1}{Your file was uploaded successfully.}
\FDISetText{12}{5}{2}{PLING!}
\FDISetText{12}{6}{1}{Now it gets\\exciting!}
\FDISetText{12}{6}{2}{End of\\part one}

% Acknowledgements
\FDISetText{13}{1}{1}{The authors would like to thank the Federal Government and the Heads of Government of the Länder, as well as the Joint Science Conference (GWK), for their funding and support within the framework of the NFDI4ING consortium. Funded by the German Research Foundation (DFG) - project number 442146713 \cite{DFG22} for NFDI4ING.}
\FDISetText{13}{1}{2}{The authors would like to thank the Federal Government and the Heads of Government of the Länder, as well as the Joint Science Conference (GWK), for their funding and support within the framework of the NFDI4ING consortium. Funded by the German Research Foundation (DFG), project number 442146713 \cite{noauthor_dfg_nodate}. The authors would further like to thank Lukas C. Bossert, Zeynep Dirier, Thomas Eifert, Marius Politze, Heinrich Schlange-Schönigen, Ines Scholten-Strauch, Benedikt Schmitz, and Annett Schwarz for discussions and helpful comments during the course of the project.}